% ============================================================================
\chapter{块设备与磁盘文件系统}
\label{ch:diskfs}
% Stage E-4 ~ E-7
% ============================================================================

\section{本章目标}

\begin{itemize}
  \item 理解块设备抽象与扇区寻址
  \item 设计可插拔 \texttt{blkdev\_ops\_t} 接口，解耦磁盘驱动
  \item 实现 ATA PIO 驱动后端
  \item 理解 ext2 磁盘布局（超级块 / 块组 / inode 表）
  \item 实现 ext2 只读文件系统后端（\texttt{fs\_ops\_t}）
  \item 通过 \texttt{KCONFIG\_BLKDEV\_BACKEND} 一键切换
\end{itemize}

% ============================================================================
\section{概念：块设备}
\label{sec:diskfs-blkdev}
% ============================================================================

% TODO:
% - 块设备 vs 字符设备：
%     块设备：固定大小块（扇区 512B / 4KB）随机读写（磁盘、SSD、USB）
%     字符设备：字节流顺序读写（串口、键盘、终端）
% - 块设备抽象：隐藏不同磁盘控制器的差异
%     ATA PIO（简单、慢，QEMU 默认）
%     ATA DMA（高效，需要内存 DMA 映射）
%     AHCI（SATA 现代接口）
%     virtio-blk（QEMU 虚拟化优化）
% - 扇区寻址：LBA（Logical Block Addressing）
%     28-bit LBA → 最大 128 GiB
%     48-bit LBA → 最大 128 PiB

% ============================================================================
\section{可插拔接口：\texttt{blkdev\_ops\_t}}
\label{sec:diskfs-blkdev-interface}
% ============================================================================

% TODO:
% - 接口设计：
%     typedef struct blkdev_ops {
%         const char* name;                              // "ata_pio", "virtio_blk"
%         int (*init)(void);                              // 探测+初始化磁盘
%         int (*read_block)(uint32_t lba, void* buf, uint32_t count);
%         int (*write_block)(uint32_t lba, const void* buf, uint32_t count);
%         uint32_t (*get_capacity)(void);                 // 总扇区数
%     } blkdev_ops_t;
% - dispatch 层 + kconfig.h: KCONFIG_BLKDEV_BACKEND 0=ata_pio

% ============================================================================
\section{后端一：ATA PIO 驱动}
\label{sec:diskfs-ata}
% ============================================================================

% TODO:
% - ATA PIO 模式寄存器（I/O ports 0x1F0–0x1F7）
% - 28-bit LBA 读扇区流程：
%     1. 选择 drive + LBA 高 4 位 → 0x1F6
%     2. 设置扇区数 → 0x1F2
%     3. 设置 LBA 低/中/高 → 0x1F3/4/5
%     4. 发送 READ SECTORS 命令 (0x20) → 0x1F7
%     5. 等待 BSY 清除 + DRQ 置位 → 从 0x1F0 读 256 个 word
% - IRQ14 中断驱动（可选，PIO 也可用轮询）
% - [WHY] ATA PIO 简单但慢（CPU 全程参与数据传输），适合学习

% ============================================================================
\section{ext2 文件系统概念}
\label{sec:diskfs-ext2-concept}
% ============================================================================

% TODO:
% - ext2 磁盘布局图：
%     Block 0: boot block（引导扇区，不用）
%     Block 1: superblock（文件系统元数据）
%     Block Group 0..N: {块组描述符, 块 bitmap, inode bitmap, inode 表, 数据块}
% - 超级块关键字段：s_inodes_count, s_blocks_count, s_block_size
% - inode 关键字段：i_mode, i_size, i_block[15]（直接/间接/双重/三重块指针）
% - 目录项：{ inode, rec_len, name_len, file_type, name[] }

% ============================================================================
\section{ext2 只读后端}
\label{sec:diskfs-ext2-ro}
% ============================================================================

% TODO:
% - 实现 fs_ops_t：
%     ext2_mount：读 superblock → 计算块组数 → 缓存关键表
%     ext2_open：路径解析（逐级 readdir） → 找到目标 inode
%     ext2_read：通过 i_block[] 定位数据块 → blkdev_read_block → 复制到 buf
%     ext2_stat：读 inode 返回 size/mode
%     ext2_readdir：解析目录数据块中的 dir_entry
% - 间接块处理（文件 > 12 blocks 时需要 indirect block pointer）
% - QEMU 磁盘镜像：qemu-img create -f raw disk.img 32M → mkfs.ext2 disk.img

% ============================================================================
\section{ext2 写入支持（进阶）}
\label{sec:diskfs-ext2-rw}
% ============================================================================

% TODO:（进阶，可跳过）
% - 块分配：查 block bitmap → 找空闲块 → 标记已用 → 更新 superblock
% - inode 分配：查 inode bitmap → 同理
% - 目录项插入：找到目录数据块末尾 → 追加 dir_entry
% - ext2_write：分配新数据块 → 更新 inode.i_block → 写数据
% - fsync：将缓存的 superblock/inode 写回磁盘

% ============================================================================
\section{验证}
\label{sec:diskfs-verify}
% ============================================================================

% TODO:
% - 测试 1：ATA PIO 读写扇区 → 写入数据 → 读回对比
% - 测试 2：挂载 ext2 → readdir("/") → 看到预置文件
% - 测试 3：open + read 一个已知内容的文件 → 内容正确
% - shell 命令：disk read <lba> / mount / ls / cat

% ============================================================================
\section{本章小结}
% ============================================================================

通过可插拔 \texttt{blkdev\_ops\_t} 和 \texttt{fs\_ops\_t} 双层抽象，
磁盘驱动与文件系统格式完全解耦——ATA PIO 可替换为 virtio-blk，ext2 可替换为 FAT32。
下一章实现字符设备框架，将串口、键盘等硬件也纳入``一切皆文件''的统一模型。
